\documentclass[11pt]{article}
\usepackage[a4paper,margin=2.5cm]{geometry}
\usepackage{amsmath,amssymb}
\usepackage{booktabs}
\usepackage{longtable}
\usepackage[hidelinks]{hyperref}
\usepackage[numbers,sort&compress]{natbib}

\title{Quantized Assimilation and Larval Feeding Rhythms\\
\large A DEB-consistent formulation with empirical grounding}
\author{}
\date{}

\begin{document}
\maketitle

This note formalizes a \textbf{bite-quantized} representation of feeding that is \textit{consistent with standard DEB flux notation}: we keep the usual continuous DEB assimilation power \(\dot p_A\) but represent it as a train of identical intake events ("bites", "sips", "mouthfuls").
\subsubsection*{Start agnostically: define bite volume and food concentration}

Let a single bite remove a \textbf{fixed volume} of food medium, denoted by \(V_{\text{bite}}\) (the \textit{bite volume}, with units \([\mathrm{cm^3\,bite^{-1}}]\)).

Let the food in the medium be described as a DEB substrate concentration \(X\) (\textit{food concentration in medium}, \([\mathrm{Cmol\,cm^{-3}}]\)).

Then the amount of food (in C-moles) per bite is

\begin{equation}
\Delta X_{\text{bite}} = X\,V_{\text{bite}}
\qquad [\mathrm{Cmol\,bite^{-1}}].
\end{equation}
\subsubsection*{Reach energy per bite via chemical potential}

Let \(\mu_X\) denote the \textit{chemical potential of food} \([\mathrm{J\,Cmol^{-1}}]\).

Then the \textbf{energy content of one bite} is

\begin{equation}
E_{\text{bite}} = \mu_X\,\Delta X_{\text{bite}}
= \mu_X\,X\,V_{\text{bite}}
\qquad [\mathrm{J\,bite^{-1}}].
\end{equation}

\subsubsection*{Define a bite frequency and derive it from DEB assimilation}

Let \(N(t)\) be the cumulative number of bites. Define the \textbf{feeding/bite frequency} as
\(\dot N_{\text{bite}}(t) \equiv \dot N(t)\) with units \([\mathrm{bite\,time^{-1}}]\) (e.g. \([\mathrm{s^{-1}}]\) or \([\mathrm{d^{-1}}]\)).

DEB gives a continuous \textbf{assimilation power} \(\dot p_A(t)\) with units \([\mathrm{J\,time^{-1}}]\).

Quantization assumption: total \textbf{ingested food energy} is the sum of bite energies per time:

\begin{equation}
\dot p_X(t) = \dot N_{\text{bite}}(t)\,E_{\text{bite}}(t).
\end{equation}

Assimilation is a fixed fraction of ingested food energy:

\begin{equation}
\dot p_A(t) = \kappa_X\,\dot p_X(t).
\end{equation}

Therefore, the bite frequency is

\begin{equation}
\boxed{
\dot N_{\text{bite}}(t)
=
\frac{\dot p_A(t)}{\kappa_X\,E_{\text{bite}}(t)}
=
\frac{\dot p_A(t)}{\kappa_X\,\mu_X\,X(t)\,V_{\text{bite}}(t)}
}
\qquad [\mathrm{bite\,time^{-1}}].
\end{equation}


\subsubsection*{Substitute a standard DEB expression for \(\dot p_A\)}

A common DEB form is

\begin{equation}
\dot p_A(t) = f(t)\,\{\dot p_{Am}\}\,S_A(t),
\end{equation}

where \(f(t)\) is the functional response (0--1), \(\{\dot p_{Am}\}\) is the max surface-area specific assimilation \([\mathrm{J\,time^{-1}\,cm^{-2}}]\), and \(S_A(t)\) is the assimilation surface area \([\mathrm{cm^2}]\).

Then

\begin{equation}
\boxed{
\dot N_{\text{bite}}(t)
=
\frac{S_A(t)}{V_{\text{bite}}(t)}\,
\frac{f(t)\,\{\dot p_{Am}\}}{\kappa_X\,\mu_X\,X(t)}
}.
\end{equation}


\subsubsection*{Break down the functional response and half-saturation}

The functional response is
\begin{equation}
f(X) = \frac{X}{K + X}, \quad 0 \le f \le 1,
\end{equation}

Substituting \(f(X)\) into the bite-frequency expression gives
\begin{equation}
\boxed{
\dot N_{\text{bite}}(t)
=
\frac{S_A(t)}{V_{\text{bite}}(t)}\,
\frac{\{\dot p_{Am}\}}{\kappa_X\,\mu_X\,(K + X(t))}
}
\label{eq:nbite_KX}
\end{equation}

The half-saturation coefficient is
\begin{equation}
K = \frac{\{\dot J_{X m}\}}{\{\dot F_m\}}.
\end{equation}

From the DEB textbook, the maximum ingestion power per area is
\begin{equation}
\{\dot p_{Xm}\} = \mu_X\,\{\dot J_{X m}\},
\end{equation}
and the corresponding maximum assimilation power is
\begin{equation}
\{\dot p_{Am}\} = \kappa_X\,\{\dot p_{Xm}\}
= \kappa_X\,\mu_X\,\{\dot J_{X m}\}.
\end{equation}
Therefore,
\begin{equation}
\{\dot J_{X m}\} = \frac{\{\dot p_{Xm}\}}{\mu_X}
= \frac{\{\dot p_{Am}\}}{\kappa_X\,\mu_X}.
\end{equation}
so
\begin{equation}
K = \frac{\{\dot J_{X m}\}}{\{\dot F_m\}}
= \frac{\{\dot p_{Am}\}}{\kappa_X\,\mu_X\,\{\dot F_m\}}.
\end{equation}

Using Eq.~(\ref{eq:nbite_KX}) and substituting \(K = \{\dot p_{Am}\}/(\kappa_X\,\mu_X\,\{\dot F_m\})\) gives
\begin{equation}
\boxed{
\dot N_{\text{bite}}(t)
=
\frac{S_A(t)}{V_{\text{bite}}(t)}\,
\frac{\{\dot p_{Am}\}}{\kappa_X\,\mu_X\left(X(t) + \frac{\{\dot p_{Am}\}}{\kappa_X\,\mu_X\,\{\dot F_m\}}\right)}
}.
\end{equation}

Simplifying the \(\kappa_X\,\mu_X\) term yields
\begin{equation}
\boxed{
\dot N_{\text{bite}}(t)
=
\frac{S_A(t)}{V_{\text{bite}}(t)}\,
\frac{\{\dot p_{Am}\}}{\kappa_X\,\mu_X\,X(t) + \frac{\{\dot p_{Am}\}}{\{\dot F_m\}}}
}.
\end{equation}

\subsubsection*{The open modeling choice: how should \(V_{\text{bite}}\) scale with size?}

At this point, the only missing ingredient is a \textbf{growth-dependent} bite volume \(V_{\text{bite}}(t)\), i.e. \(V_{\text{bite}}(L)\). A general family with \(\kappa_{bite}\) the coefficient setting the bite magnitude per bite and \(k\) the bite-volume scaling exponent:

\begin{equation}
V_{\text{bite}}(L) = \kappa_{bite}\,L^{k}
\qquad [\mathrm{cm^3\,bite^{-1}}],
\end{equation}


Two principled options correspond to \textbf{area vs volume} scaling, with specific choices of the exponent:

\paragraph{(A) Volume-scaled bites: \(k=3\)}
\begin{equation}
V_{\text{bite}}(L)=\kappa_{bite,V}\,V(L),
\end{equation}
where \(V(L)\) is structural volume (for an isomorph, typically \(V(L)\propto L^3\)).

\textbf{Consequence (for fixed \(X\), \(f\)):}
\begin{equation}
\dot N_{\text{bite}} \propto \frac{L^2}{L^3} \propto L^{-1}.
\end{equation}
Larger juveniles take fewer bites per time, but each bite is larger.

\emph{Interpretation: if you want each bite to represent a roughly fixed fraction of body (structural) volume, choose this volume-scaling option with prefactor \(\kappa_{bite,V}\) and exponent \(k=3\).}

\paragraph{(B) Area-scaled bites: \(k=2\)}
\begin{equation}
V_{\text{bite}}(L)=\kappa_{bite,S}\,L^2.
\end{equation}

\textbf{Consequence (for fixed \(X\), \(f\)):}
\begin{equation}
\dot N_{\text{bite}} \propto \frac{L^2}{L^2} \propto L^{0},
\end{equation}
so bite frequency tends to remain approximately constant over growth (a ``feeding rhythm'' interpretation).
\emph{Interpretation: if you want a roughly constant feeding rhythm across growth, choose this area-scaling option with prefactor \(\kappa_{bite,S}\) and exponent \(k=2\).}
\newpage
\section*{Mouth-hook movement rate across larval development}

This section synthesizes \textbf{modern neuroethological / circuit-based} and \textbf{classical behavioral genetics / developmental} studies relevant to modeling bite / mouth-hook frequency across \textit{Drosophila melanogaster} larval development. Together they delimit what is known, what is measured, and what remains open for DEB-style quantized feeding models.

\paragraph{Modern neuroethological \& circuit-based literature}

\paragraph{Scope}
Recent studies focus on the \textbf{larval feeding circuit} (brain + subesophageal zone), use \textbf{mouth-hook contractions} as a measurable behavioral output, and examine modulation by \textbf{hunger, neuromodulators, and developmental state} (see also the review by \citet{miroschnikow2020curbio}). These studies usually work with \textbf{late L2 or L3 larvae}, because circuits are mature and behavior is robust.

\paragraph{Key findings}
Mouth-hook contraction rate emerges as a \textbf{stable, high-frequency rhythmic behavior during active feeding}. Typical reported values cluster around \textbf{150--180 contractions/min} during active feeding on yeast or standard media, and feeding rate \textbf{drops sharply} during the wandering/pre-pupal transition. Developmental comparisons are \textbf{rarely systematic across all instars}, but authors note that rates are lower and more variable in early instars, relatively stable during most of L3, and actively suppressed near pupation.

\begin{table}[h]
\centering
\begin{tabular}{llll}
\toprule
Stage & Reported rate & Notes & Source \\
\midrule
late L2 / early L3 & $\sim 160\text{--}170\,\text{min}^{-1}$ & Active feeding on yeast & \citep{bhatt2013jove} \\
mid L3 & $\sim 150\text{--}180\,\text{min}^{-1}$ & Depends on hunger state & \citep{schoofs2018jinsphys} \\
late L3 & $\downarrow$ strongly & Feeding actively suppressed & \citep{melcher2005plosbio} \\
\bottomrule
\end{tabular}
\\[0.5em]
Table 1: Quantitative anchor values for mouth-hook contraction rates in feeding assays.
\end{table}
\paragraph{Classical behavioral genetics \& developmental feeding studies}

Earlier quantitative work explicitly examined \textbf{feeding rate vs larval age}, genetic variation in feeding behavior (e.g. rover/sitter phenotypes), and mouth-hook or cephalopharyngeal retraction counts as a \textit{direct metric}.

\begin{table}[h]
\centering
\begin{tabular}{llll}
\toprule
Developmental phase & Feeding rate trend & Interpretation & Source \\
\midrule
Early larva (L1) & Low, rising & Small mouthparts, immature circuit & \citep{sewell1974genetres,burnet1977genetres} \\
L2 & Increasing & Growth acceleration & \citep{sewell1974genetres,burnet1977genetres} \\
Mid L3 & Plateau & Sustained biomass accumulation & \citep{sewell1974genetres,burnet1977genetres} \\
Wandering / pre-pupa & Decline & Feeding inhibition precedes metamorphosis & \citep{sewell1974genetres,burnet1977genetres,debelle1987heredity} \\
\bottomrule
\end{tabular}
\\[0.5em]
Table 2: Developmental changes in larval feeding rate and interpretation.
\end{table}
\subsubsection*{Synthesis, implications, and bottom line}

The combined modern and classical literatures agree that mouth-hook contraction rate is a \textbf{valid and widely used proxy} for feeding/bite frequency. Rates are \textbf{developmentally regulated}, not constant across instars: they rise early, plateau during much of L3, and are \textbf{actively suppressed} near pupation, with a \textbf{plateau during most of L3} emerging as a particularly strong empirical feature across studies.

Key gaps remain, including the lack of a \textbf{single standardized dataset} measuring mouth-hook frequency continuously from L1 -> L2 -> L3 under identical conditions, and the absence of direct, joint measurement of bite \textbf{frequency}, bite \textbf{duration}, and bite \textbf{volume} within the same experiments.

For DEB-style bite quantization, these findings imply that \textbf{constant bite frequency across the full larval period is empirically false}, although within L3 a roughly constant frequency is a reasonable approximation. Growth-related changes in assimilation are therefore most naturally attributed to changes in \textbf{bite volume} and in the \textbf{fraction of time spent feeding}, rather than to unbounded increases in frequency. Mouth-hook frequency can be treated as an \textit{upper bound / rhythm}, with effective feeding rate given by frequency $\times$ bite volume $\times$ feeding duty cycle.

Taken together, these points make bite-quantized DEB models \textit{biologically plausible} and highlight exactly where new data (continuous developmental time series and joint kinematic measurements) would be most informative.

\appendix
\section*{Appendix: Rod-shape isomorph choice (shape correction included)}
For your \textbf{rod-shaped isomorph} with a \textbf{shape correction} \(\delta_M\), a convenient mapping is to retain DEB's length-based scaling but correct for shape. The structural volume is
\begin{equation}
V(t) = (\delta_M L(t))^3
\qquad [\mathrm{cm^3}]
\end{equation}

and the assimilation area (surface scaling) is
\begin{equation}
S_A(t) = (\delta_M L(t))^2
\qquad [\mathrm{cm^2}]
\end{equation}

and thus

\begin{equation}
\dot p_A(t) = f(t)\,\{\dot p_{Am}\}\,(\delta_M L(t))^2.
\end{equation}

So the bite frequency becomes

\begin{equation}
\boxed{
\dot N_{\text{bite}}(t)
=
\frac{(\delta_M L(t))^2}{V_{\text{bite}}(t)}\,
\frac{f(t)\,\{\dot p_{Am}\}}{\kappa_X\,\mu_X\,X(t)}
}.
\end{equation}

\begin{quote}
Note: if you later choose to use an explicit rod lateral area formula (rather than the \((\delta_M L)^2\) isomorph mapping), you can swap \(S_A(t)\) accordingly without changing the rest of the derivation.
\end{quote}


\bibliographystyle{plainnat}
\bibliography{quantized_assimilation_and_feeding}

\end{document}
